\documentclass[12pt]{article}
\usepackage{amsmath, amssymb}
\usepackage{geometry}
\usepackage{booktabs}
\usepackage{hyperref}

\geometry{margin=1in}

\title{MSE 435 -- Progress Report 2\\[0.5em]
\large Examination Room Scheduling in an Outpatient Clinic}
\author{
  Johnny Cao: 20949892 \\
  Chintan Desai: 20938455 \\
  Anastan Gnanapragasam: 20948196
}
\date{MSE 435 \\ Winter 2026}

\begin{document}

\maketitle

% ─────────────────────────────────────────────
\section{Problem Overview}

The goal of this project is to assign appointments to examination rooms and start times in a way
that respects operational constraints of the clinic while minimizing unnecessary room switching by
healthcare providers. The clinic operates within fixed hours and includes administrative blocks and
a lunch closure. To model this accurately, time is discretized into 15-minute intervals. This provides
enough granularity to represent appointment durations while keeping the model manageable. We
first formulate a time-indexed mixed-integer program. Since this model can grow large, we then
reformulate it using column generation.

% ─────────────────────────────────────────────
\section{Time-Indexed Assignment Formulation}

\subsection*{Sets}
\begin{itemize}
  \item $I$ = set of appointments
  \item $P$ = set of providers
  \item $R$ = set of examination rooms
  \item $T$ = set of 15-minute time slots
\end{itemize}

Each appointment $i \in I$ has duration $d_i$ (in time slots) and a designated provider $p(i)$.

\subsection*{Parameters}
\begin{itemize}
  \item $A_{pt} = 1$ if provider $p$ is available at time $t$
  \item $C_{ir} = 1$ if appointment $i$ can use room $r$
  \item $\delta_{rr'} =$ distance between rooms $r$ and $r'$
\end{itemize}

\subsection*{Decision Variable}
\[
  x_{itr} =
  \begin{cases}
    1 & \text{if appointment } i \text{ starts at time } t \text{ in room } r \\
    0 & \text{otherwise}
  \end{cases}
\]

\subsection*{Objective Function}

The main operational objective is to reduce provider movement between rooms. We model total
travel distance between consecutive appointments for each provider. This encourages compact room
usage and discourages unnecessary switching. Introducing auxiliary variables $y_{ii'rr'}$ to represent
direct transitions between appointments, we minimize:
\[
  \min \sum_{p \in P} \sum_{\substack{i,i' \in I \\ p(i)=p(i')=p}} \sum_{r,r' \in R} \delta_{rr'}\, y_{ii'rr'}.
\]

We define
\[
  y_{ii'rr'} =
  \begin{cases}
    1 & \text{if appointment } i \text{ is immediately followed by appointment } i' \text{ for the same provider,} \\
      & \text{with } i \text{ in room } r \text{ and } i' \text{ in room } r' \\
    0 & \text{otherwise.}
  \end{cases}
\]

In this progress report, we treat these transition indicators as being derived from the chosen
appointment start-time assignments; in implementation, they are computed from the final schedule
produced by the model (or enforced via additional sequencing constraints when needed).

\subsection*{Constraints}

\paragraph{(1) Each appointment scheduled exactly once.}
This ensures full coverage.
\[
  \sum_{t \in T} \sum_{r \in R} x_{itr} = 1 \qquad \forall i \in I.
\]

\paragraph{(2) Room capacity constraints.}
A room cannot host overlapping appointments:
\[
  \sum_{i \in I} \sum_{\tau \in T:\, \tau \le t < \tau + d_i} x_{i\tau r} \le 1
  \qquad \forall r \in R,\, t \in T.
\]

\paragraph{(3) Provider capacity constraints.}
A provider cannot see more than one patient at a time:
\[
  \sum_{i:\, p(i)=p} \sum_{r \in R} \sum_{\tau \in T:\, \tau \le t < \tau + d_i} x_{i\tau r} \le 1
  \qquad \forall p \in P,\, t \in T.
\]

\paragraph{(4) Provider availability.}
Appointments must lie fully within provider working hours:
\[
  x_{itr} = 0 \quad \text{if } \{t, \ldots, t + d_i - 1\} \not\subseteq A_{p(i)}.
\]
Equivalently, we enforce:
\[
  x_{itr} \le A_{p(i),t'} \qquad \forall i \in I,\, \forall r \in R,\, \forall t \in T,\, \forall t' \in \{t, \ldots, t + d_i - 1\}.
\]

\paragraph{(5) Room compatibility.}
These constraints define the feasible region of the model.
\[
  x_{itr} = 0 \quad \text{if } C_{ir} = 0.
\]
\[
  x_{itr} \le C_{ir} \qquad \forall i \in I,\, \forall t \in T,\, \forall r \in R.
\]

% ─────────────────────────────────────────────
\section{Column Generation Reformulation}

Although the above formulation is intuitive, the number of variables is on the order of
$|I||T||R|$, which becomes large for weekly planning. To improve scalability, we reformulate
the problem using Dantzig--Wolfe decomposition.

\subsection*{Column Definition}

For each provider-day block $g$, let $K_g$ denote the set of all feasible daily schedules for
that block. Each column $k \in K_g$ represents:
\begin{itemize}
  \item a feasible ordering of that block's appointments
  \item assigned start times and rooms
  \item no overlaps
  \item compliance with availability and compatibility constraints
\end{itemize}

A provider-day block is the basic unit in the implementation. Because the raw appointment
data contains same-provider overlaps on some days, each original provider-day is refined into
one or more serial blocks before column generation. This is a practical refinement in Part 3
and does not change the master problem structure.

\subsection*{Column Cost}

If all appointments occur in the same room, the cost is zero. The cost of a schedule $k$ is
the total travel distance between rooms in that schedule:
\[
  c_k = \sum_{\text{room transitions within schedule } k} \delta_{rr'}.
\]

\subsection*{Set Partitioning}

We now choose exactly one schedule per provider-day block.

\paragraph{Objective Function.}
\[
  \min \sum_{g} \sum_{k \in K_g} c_k \lambda_k
\]

\paragraph{Subject to:}

\noindent\textit{Appointment coverage}
\[
  \sum_{g} \sum_{k \in K_g} a_{ik} \lambda_k = 1 \qquad \forall i \in I.
\]

\noindent\textit{One schedule per provider-day block}
\[
  \sum_{k \in K_g} \lambda_k = 1 \qquad \forall g.
\]

\noindent\textit{Clinic-wide room capacity}
\[
  \sum_{g} \sum_{k \in K_g} b_{krt} \lambda_k \le 1 \qquad \forall r \in R,\, t \in T.
\]
\[
  \lambda_k \in \{0, 1\}.
\]

In practice, $K_g$ is very large, so we solve a restricted master problem (RMP) using a small
subset of columns and add new columns via pricing.

% ─────────────────────────────────────────────
\section{Solution Methodology: Column Generation Approach}

To solve large instances of the scheduling problem, we use column generation based on the
set-partitioning reformulation above. The idea is to begin with a restricted master problem
containing only a small number of feasible schedules for each provider-day block, and then
iteratively generate new schedules only when they can improve the objective.

\subsection*{Step 1: Initialize a Restricted Master Problem}

We begin with a small subset of columns for each block, denoted $K^0_g \subset K_g$, such
that every appointment is covered by at least one initial schedule. These initial columns are
constructed using simple feasible patterns, such as assigning each block's appointments in
chronological order to currently allowable rooms.

The initial restricted master problem is:
\[
  \min \sum_{g} \sum_{k \in K^0_g} c_k \lambda_k
\]
subject to
\[
  \sum_{g} \sum_{k \in K^0_g} a_{ik} \lambda_k = 1 \qquad \forall i \in I,
\]
\[
  \sum_{k \in K^0_g} \lambda_k = 1 \qquad \forall g,
\]
\[
  \sum_{g} \sum_{k \in K^0_g} b_{krt} \lambda_k \le 1 \qquad \forall r \in R,\, t \in T,
\]
\[
  \lambda_k \ge 0.
\]
At this stage, we solve the LP relaxation of the master problem.

\subsection*{Step 2: Solve the Restricted Master Problem}

Solving the LP relaxation of the RMP gives dual values associated with the master constraints. Let:
\begin{itemize}
  \item $\pi_i$ be the dual variable for appointment coverage constraint $i$
  \item $\sigma_g$ be the dual variable for the one-schedule-per-block constraint for block $g$
  \item $\mu_{rt}$ be the dual variable for the room-capacity constraint for room $r$ and time slot $t$
\end{itemize}

These dual values are then used to define the reduced cost of a potential new schedule.

\subsection*{Step 3: Solve Provider-Day Pricing Subproblems}

For each provider-day block $g$, we solve a pricing subproblem to search for a new feasible
schedule with negative reduced cost. The reduced cost of a schedule $k \in K_g$ is:
\[
  \bar{c}_k = c_k - \sum_{i \in I} \pi_i a_{ik} - \sigma_g + \sum_{r \in R} \sum_{t \in T} \mu_{rt} b_{krt}.
\]

Thus, for each block $g$, the pricing problem is:
\[
  \min_{k \in K_g} \left( c_k - \sum_{i \in I} \pi_i a_{ik} - \sigma_g + \sum_{r \in R} \sum_{t \in T} \mu_{rt} b_{krt} \right).
\]

Each pricing subproblem generates a feasible schedule that respects appointment durations,
provider availability, room compatibility, and no-overlap requirements. If the minimum reduced
cost is negative, then that schedule can improve the current master solution and is added to the RMP.

\subsection*{Step 4: Add Improving Columns and Repeat}

After solving the pricing subproblems for all provider-day blocks, any schedule with negative
reduced cost is added to the restricted master problem. The updated RMP is then re-solved,
producing a new set of dual values. This process is repeated until no block can generate a
schedule with negative reduced cost.

At that point, the LP relaxation of the full master problem has been solved.

\subsection*{Step 5: Recover an Integer Solution}

The solution obtained from column generation is generally fractional because the restricted
master problem is solved as an LP relaxation. To obtain a feasible scheduling plan, we then
enforce integrality on the master variables:
\[
  \lambda_k \in \{0, 1\}.
\]
An exact approach is to apply branch-and-bound or branch-and-cut on the master problem
after column generation has converged. In the implementation, the final generated column pool
is fixed and the resulting restricted integer master problem is solved directly.

\subsection*{Step 6: Stopping Criteria}

The algorithm ends when:
\begin{itemize}
  \item no pricing subproblem generates a column with negative reduced cost, and
  \item the resulting restricted integer master problem has been solved to feasibility or optimality.
\end{itemize}

The final schedule is the best feasible solution obtained from the selected provider-day schedules.

\subsection*{Remark}

Room preference is tracked as a practical metric in the implementation, but the optimization
objective minimized by the master problem is room-switch distance only. Preference penalties
are reported separately in the results as a secondary schedule-quality metric.

% ─────────────────────────────────────────────
\section{Solving the Model and Practical Refinements}

We solve the resulting model in two stages. First, the time-indexed formulation is solved on
reduced daily instances to verify that the feasibility constraints and objective behave as intended.
Second, the full weekly problem is solved using the column generation framework, with additional
refinements to ensure that the selected schedules are practical for clinic use.

\subsection*{Data Preprocessing and Feasibility Filtering}

The appointment data is first converted into the discrete time-slot representation used in the
model. Appointment durations are rounded up to the nearest 15-minute interval. Cancelled and
deleted appointments are removed from the active demand set, while no-show appointments are
retained because they still consume capacity at the planning stage.

Next, infeasible appointment placements are filtered out before optimization. Any start time
that causes an appointment to overlap with the clinic lunch closure or administrative blocks is
removed from the feasible placement set when generating columns. In addition, the provider-room
assignment tables are converted into time-dependent feasibility rules so that AM/PM room
restrictions are respected when building feasible schedules. This preprocessing step reduces the
size of the feasible region and ensures that the optimization only considers implementable assignments.

Because the raw data contains overlapping appointments for the same provider on the same
day, each original provider-day is refined into the minimum number of non-overlapping serial
blocks before solving. Column generation is then carried out on these refined serial blocks. This
is the main practical refinement used to make the model consistent with the real data.

\subsection*{Baseline Validation on Reduced Instances}

Before solving the full problem, the time-indexed model is used on reduced instances as a
validation model. This serves two purposes. First, it confirms that the coverage, provider
capacity, room capacity, availability, and compatibility constraints correctly produce feasible
schedules. Second, it provides a benchmark that can be compared against the column generation
solution on smaller blocks.

In the implementation, small refined blocks are solved by both the validation model and the
pricing-based schedule generator, and the solutions are compared. The final results show zero
validation mismatches in both weeks.

\subsection*{Initial Column Generation}

To initialize the restricted master problem, we generate a feasible starting schedule for each
refined provider-day block in the planning horizon. Appointments are ordered chronologically,
then assigned to the earliest feasible start time and the most preferred allowable room. If the
preferred room is unavailable, the nearest compatible room is chosen using the room distance matrix.

This produces an initial set of columns that is feasible and operationally reasonable, while
also giving the master problem a meaningful starting point.

\subsection*{Pricing Subproblem Implementation}

In the implementation, pricing is solved separately for each refined provider-day block. To
produce meaningful schedules and keep the subproblem tractable, the chronological order of
appointments within each block is preserved, while room assignments are optimized around that order.

Let the appointments for a given block be indexed in chronological order as $i_1, \ldots, i_m$.
For each appointment $i_q$, we consider only feasible placements $(t, r)$ that satisfy duration,
provider availability, room compatibility, and clinic operating rules. The pricing problem is then
solved by dynamic programming.

Define $F_q(t, r)$ as the minimum reduced cost of scheduling the first $q$ appointments, where
appointment $i_q$ starts at time $t$ in room $r$. For the first appointment:
\[
  F_1(t, r) = -\pi_{i_1} - \sigma_g + \sum_{u=t}^{t+d_{i_1}-1} \mu_{ru}.
\]

For $q = 2, \ldots, m$, the recurrence is
\[
  F_q(t, r) = -\pi_{i_q} + \sum_{u=t}^{t+d_{i_q}-1} \mu_{ru}
  + \min_{(s,r') \prec (t,r)} \left\{ F_{q-1}(s, r') + \delta_{r'r} \right\},
\]
where $(s, r') \prec (t, r)$ means that appointment $i_{q-1}$ placed at $(s, r')$ finishes before
appointment $i_q$ starts at $t$, and both placements are feasible.

The reduced cost of the best full schedule for block $g$ is then
\[
  \bar{c}_g = \min_{t,r} F_m(t, r).
\]
If $\bar{c}_g < 0$, the corresponding schedule is added to the restricted master problem as a new
column. This process is repeated until no pricing problem can produce a schedule with negative
reduced cost.

\subsection*{Final Integer Solution}

After column generation terminates, the final column pool is fixed and the restricted master
problem is re-solved with integrality enforced:
\[
  \lambda_k \in \{0, 1\}.
\]
This restricted integer master problem selects the final block schedules while respecting
clinic-wide room capacity. Solving the integer master on the generated column pool gives a
complete feasible weekly schedule.

\subsection*{Practical Refinements for Meaningful Schedules}

Several refinements are used so that the final schedules are useful in practice rather than only
mathematically feasible.

First, the pricing routine preserves each block's within-day appointment order. This avoids
unrealistic patient reordering and makes the generated schedules easier to interpret operationally.

Second, when multiple columns have similar reduced cost, preference is given to schedules with
fewer room switches and stronger consistency with the provider's preferred room assignment.
This keeps providers in the same room whenever possible and reduces unnecessary movement.

Third, if the final restricted integer master problem is infeasible or produces overly fragmented
room usage, additional columns are generated by relaxing room preference while keeping all
hard feasibility constraints unchanged. This provides a controlled refinement mechanism without
changing the underlying model structure.

\subsection*{Output Schedule}

The final output of the solution procedure is a feasible weekly clinic schedule specifying, for
each appointment, its assigned provider, start time, end time, and examination room. These
schedules can then be summarized using provider-day tables and Gantt-style visualizations for
managerial interpretation in the next stage of the project.

% ─────────────────────────────────────────────
\section{Computational Results}

\subsection*{Instance Characteristics}

After removing cancelled and deleted appointments, Week 1 contained 427 active appointments
across 20 providers and 5 clinic days, while Week 2 contained 499 active appointments across 23
providers and 5 clinic days. The raw booked data was not directly serial by provider: Week 1
contained 73 same-provider overlaps in the input, while Week 2 contained 110. Similarly, 44
appointments in Week 1 and 58 appointments in Week 2 touched administrative or lunch blocks
in the raw data, which justified the feasibility filtering step in the implementation.

At the provider-day level, Week 1 contained 64 original provider-day blocks that were refined
into 94 serial blocks, while Week 2 contained 79 original provider-day blocks that were refined
into 122 serial blocks. This confirms that the refinement step was necessary to make the
column-generation model workable on the real appointment files.

\subsection*{Duration and Daily Demand Structure}

The demand in both weeks was dominated by short appointments. For Week 1, the durations
converted into slot counts as follows: 224 one-slot appointments, 182 two-slot appointments,
8 three-slot appointments, 11 four-slot appointments, 1 six-slot appointment, and 1 eight-slot
appointment. For Week 2, the durations converted into 296 one-slot appointments, 171 two-slot
appointments, 12 three-slot appointments, and 20 four-slot appointments.

Daily demand was also uneven. Week 1 active appointments by day were 125, 4, 108, 120,
and 70 from Monday to Friday. Week 2 active appointments by day were 141, 132, 104, 42, and
80. These patterns support solving the model by refined provider-day blocks rather than as one
monolithic weekly model.

\subsection*{Column Generation Performance}

Table~\ref{tab:results} summarizes the main computational results.

The column pool grew meaningfully in both weeks, from 94 to 160 columns in Week 1 and from
122 to 247 columns in Week 2. The method converged in 8 iterations for Week 1 and 13
iterations for Week 2. In both cases, the final LP objective matched the final integer objective,
indicating that the final restricted master solution was already integral at convergence.

\subsection*{Schedule Quality Metrics}

The implemented optimization objective minimizes room-switch distance only. For that reason,
the final integer objective should be interpreted as the optimized room-switch cost from the
restricted master problem. Room preference penalties were tracked separately as secondary
metrics and were not part of the master objective.

The final reported schedule metrics were:
\begin{itemize}
  \item Week 1: final LP objective $= 51.9$; final integer objective $= 51.9$; room-switch distance $= 339.7$; preference penalty $= 318.7$; runtime $= 176.87$ s
  \item Week 2: final LP objective $= 118.8$; final integer objective $= 118.8$; room-switch distance $= 600.9$; preference penalty $= 525.3$; runtime $= 311.39$ s
\end{itemize}

These values are larger than the final master objective because they are computed after schedule
reconstruction and include post-processed schedule-quality measures rather than only the
switch-only objective optimized by the master problem. They are therefore reported separately
and should not be added directly to the final integer objective.

\begin{table}[h]
  \centering
  \caption{Column generation results by week.}
  \label{tab:results}
  \begin{tabular}{lrr}
    \toprule
    Metric & Week 1 & Week 2 \\
    \midrule
    Active appointments              & 427    & 499    \\
    Providers                        & 20     & 23     \\
    Original provider-day blocks     & 64     & 79     \\
    Refined serial blocks            & 94     & 122    \\
    Initial columns                  & 94     & 122    \\
    Final column pool size           & 160    & 247    \\
    Column generation iterations     & 8      & 13     \\
    Final LP objective               & 51.9   & 118.8  \\
    Final integer objective          & 51.9   & 118.8  \\
    Runtime (seconds)                & 176.87 & 311.39 \\
    Validated small blocks           & 74     & 99     \\
    Validation mismatches            & 0      & 0      \\
    \bottomrule
  \end{tabular}
\end{table}

\subsection*{Interpretation}

The results show that the refined column-generation approach was able to solve both weekly
instances and produce feasible schedules with zero validation mismatches on the tested small
blocks. The implementation also handled the main practical complications in the data, namely
overlapping same-provider appointments and partial-day room restrictions, by refining
provider-day blocks before optimization and by filtering infeasible placements during schedule generation.

Overall, the computational results support the choice of a decomposed set-partitioning model
over a direct full-week time-indexed formulation. The number of active appointments, the
uneven daily demand, and the need to enforce clinic-wide room capacity all make the large-scale
reformulation more appropriate for the actual clinic scheduling data.

\end{document}